%% LaTeX Commands and Macros for Lambdust Language Specification
%% This file defines specialized typography and formatting for the specification
%% Includes R7RS-style commands for compatibility

\makeatletter

% R7RS chapter/section commands
\newcommand{\topnewpage}{\@topnewpage}

\newcommand{\clearchapterstar}[1]{
  \clearpage
  \topnewpage[
    \centerline{\large\bf\uppercase{#1}}
    \bigskip]}

% R7RS formal semantics notation
\def\elem{\hbox{\raise.13ex\hbox{$\scriptstyle\in$}}}

\makeatother

% Font setup for Unicode code support
% Try to use JuliaMono if available, fallback to system monospace
\IfFontExistsTF{JuliaMono-Regular.ttf}{
    \setmonofont{JuliaMono-Regular.ttf}[Scale=0.9]
}{
    \IfFontExistsTF{JuliaMono}{
        \setmonofont{JuliaMono}[Scale=0.9]
    }{
        \setmonofont{Menlo}[Scale=0.9]  % macOS fallback
    }
}

% Unicode character definitions for common symbols in Lambdust code
\newunicodechar{λ}{\ensuremath{\lambda}}
\newunicodechar{π}{\ensuremath{\pi}}
\newunicodechar{Π}{\ensuremath{\Pi}}
\newunicodechar{Σ}{\ensuremath{\Sigma}}
\newunicodechar{→}{\ensuremath{\rightarrow}}
\newunicodechar{⊢}{\ensuremath{\vdash}}
\newunicodechar{∀}{\ensuremath{\forall}}
\newunicodechar{∃}{\ensuremath{\exists}}
\newunicodechar{≤}{\ensuremath{\leq}}
\newunicodechar{≥}{\ensuremath{\geq}}
\newunicodechar{≠}{\ensuremath{\neq}}
\newunicodechar{∈}{\ensuremath{\in}}
\newunicodechar{∉}{\ensuremath{\notin}}
\newunicodechar{∅}{\ensuremath{\emptyset}}
\newunicodechar{∪}{\ensuremath{\cup}}
\newunicodechar{∩}{\ensuremath{\cap}}
\newunicodechar{⊆}{\ensuremath{\subseteq}}
\newunicodechar{⊇}{\ensuremath{\supseteq}}
\newunicodechar{⟨}{\ensuremath{\langle}}
\newunicodechar{⟩}{\ensuremath{\rangle}}

% Code listing style
\definecolor{lambdust-comment}{rgb}{0.5,0.5,0.5}
\definecolor{lambdust-keyword}{rgb}{0.0,0.0,1.0}
\definecolor{lambdust-string}{rgb}{0.0,0.6,0.0}
\definecolor{lambdust-number}{rgb}{0.8,0.0,0.8}
\definecolor{lambdust-type}{rgb}{0.6,0.0,0.6}

\lstdefinestyle{lambdust}{
    language=Lisp,
    basicstyle=\ttfamily\footnotesize,
    commentstyle=\color{lambdust-comment}\itshape,
    keywordstyle=\color{lambdust-keyword}\bfseries,
    stringstyle=\color{lambdust-string},
    numberstyle=\color{lambdust-number},
    showstringspaces=false,
    breaklines=true,
    breakatwhitespace=true,
    tabsize=2,
    frame=single,
    rulecolor=\color{gray!30},
    backgroundcolor=\color{gray!5},
    captionpos=b,
    % Lambdust-specific keywords
    morekeywords={define,lambda,let,let*,letrec,if,cond,case,and,or,not,
                  begin,do,when,unless,quote,quasiquote,unquote,syntax,
                  syntax-case,syntax-rules,with-syntax,datum,
                  Pi,Sigma,Refinement,the,perform,handle,lift,
                  spawn,send,receive,future,await,
                  foreign-call,foreign-data,capability,
                  prove,extract,qed,
                  Dynamic,Number,String,Boolean,Symbol,List,Vector,
                  Effect,Actor,Future,Capability},
    % Type annotations
    moredelim=[s][\color{lambdust-type}]{\#:}{ },
    moredelim=[s][\color{lambdust-type}]{::}{ },
    moredelim=[s][\color{lambdust-type}]{->}{ },
    moredelim=[s][\color{lambdust-type}]{=>}{ },
    moredelim=[s][\color{lambdust-type}]{\~>}{ },
}

% Inline code style
\newcommand{\code}[1]{\texttt{\small #1}}
\newcommand{\scheme}[1]{\lstinline[style=lambdust]!#1!}

% Semantic notation (from formal semantics)
\newcommand{\sembrack}[1]{[\![#1]\!]}
\newcommand{\fun}[1]{\hbox{\it #1}}
\newenvironment{semfun}{\begin{tabbing}$}{$\end{tabbing}}

% Domain notation
\newcommand\LOC{{\tt{}L}}
\newcommand\NAT{{\tt{}N}}
\newcommand\TRU{{\tt{}T}}
\newcommand\SYM{{\tt{}Q}}
\newcommand\CHR{{\tt{}H}}
\newcommand\NUM{{\tt{}R}}
\newcommand\FUN{{\tt{}F}}
\newcommand\EXP{{\tt{}E}}
\newcommand\STV{{\tt{}E}}
\newcommand\STO{{\tt{}S}}
\newcommand\ENV{{\tt{}U}}
\newcommand\ANS{{\tt{}A}}
\newcommand\ERR{{\tt{}X}}
\newcommand\DP{\tt{P}}
\newcommand\EC{{\tt{}K}}
\newcommand\CC{{\tt{}C}}
\newcommand\MSC{{\tt{}M}}
\newcommand\PAI{\hbox{\EXP$_{\rm p}$}}
\newcommand\VEC{\hbox{\EXP$_{\rm v}$}}
\newcommand\STR{\hbox{\EXP$_{\rm s}$}}

% Lambdust extension domains
\newcommand\TYP{{\tt{}T}}     % Types
\newcommand\EFF{{\tt{}F}}     % Effects  
\newcommand\MON{{\tt{}M}}     % Monads
\newcommand\ACT{{\tt{}A}}     % Actors
\newcommand\CON{{\tt{}C}}     % Concurrency primitives
\newcommand\FFI{{\tt{}I}}     % FFI bindings
\newcommand\DEP{{\tt{}D}}     % Dependent types
\newcommand\PRF{{\tt{}P}}     % Proof terms

% Grammar notation
\newcommand{\goesto}{$\longrightarrow$}
\newcommand{\produces}{\ensuremath{\to}}
\newcommand{\arbno}[1]{{#1}$^*$}
\newcommand{\atleastone}[1]{{#1}$^+$}
\newcommand{\opt}[1]{{#1}$^?$}
\newcommand{\alternative}{\ensuremath{\ |\ }}

% Grammar environment
\newenvironment{grammar}{
    \begin{center}
    \begin{tabular}{r@{\quad}c@{\quad}l@{\qquad}l}
}{
    \end{tabular}
    \end{center}
}

% BNF production rules
\newcommand{\production}[4]{
    #1 & #2 & #3 & #4 \\
}

% Auxiliary notation
\newcommand\elt{\downarrow}
\newcommand\drop{\dagger}
\newcommand{\wrong}[1]{\fun{wrong }\hbox{\rm``#1''}}
\newcommand{\go}[1]{\hbox{\hspace*{#1em}}}

% Type system notation
\newcommand{\typeof}[1]{\ensuremath{\tau(#1)}}
\newcommand{\hastype}[2]{\ensuremath{#1 : #2}}
\newcommand{\typecheck}[3]{\ensuremath{#1 \vdash #2 : #3}}
\newcommand{\subtype}[2]{\ensuremath{#1 \leq #2}}
\newcommand{\consistent}[2]{\ensuremath{#1 \sim #2}}
\newcommand{\gradual}[1]{\ensuremath{\mathcal{G}(#1)}}

% Effect notation
\newcommand{\lambdusteffects}[1]{\ensuremath{\varepsilon(#1)}}
\newcommand{\lambdustpure}{\ensuremath{\emptyset}}
\newcommand{\lambdusteffectful}[1]{\ensuremath{\{#1\}}}
\newcommand{\lambdustunion}{\cup}
\newcommand{\effectarrow}[2]{\ensuremath{#1 \rightsquigarrow #2}}

% Concurrency notation
\newcommand{\lambdustspawn}[1]{\ensuremath{\text{spawn}(#1)}}
\newcommand{\lambdustsend}[2]{\ensuremath{#1 ! #2}}
\newcommand{\receive}{\ensuremath{?}}
\newcommand{\lambdustparallel}[2]{\ensuremath{#1 \parallel #2}}

% Proof notation
\newcommand{\proves}[2]{\ensuremath{#1 \vdash #2}}
\newcommand{\lambdustextract}[1]{\ensuremath{\text{extract}(#1)}}
\newcommand{\lambdustqed}{\ensuremath{\square}}

% Special forms
\newcommand{\lambdaform}[2]{\ensuremath{(\lambda\,#1\,.\,#2)}}
\newcommand{\letform}[3]{\ensuremath{(\text{let}\,([#1\,#2])\,#3)}}
\newcommand{\ifform}[3]{\ensuremath{(\text{if}\,#1\,#2\,#3)}}

% Syntax highlighting for special symbols
\newcommand{\lambdustkeyword}[1]{\textcolor{lambdust-keyword}{\textbf{#1}}}
\newcommand{\lambdusttype}[1]{\textcolor{lambdust-type}{\textsf{#1}}}
\newcommand{\lambdusteffect}[1]{\textcolor{red}{\textit{#1}}}

% Box environments
\newenvironment{syntaxbox}{
    \begin{center}
    \fbox{\begin{minipage}{0.9\textwidth}
    \vspace{0.5em}
}{
    \vspace{0.5em}
    \end{minipage}}
    \end{center}
}

\newenvironment{examplebox}{
    \begin{center}
    \colorbox{gray!10}{\begin{minipage}{0.9\textwidth}
    \vspace{0.5em}
}{
    \vspace{0.5em}
    \end{minipage}}
    \end{center}
}

% Custom theorem style for examples with line breaks
\newtheoremstyle{examplestyle}%  name
  {\topsep}%                     space above
  {\topsep}%                     space below  
  {}%                            body font
  {}%                            indent amount
  {\bfseries}%                   theorem head font
  {.}%                           punctuation after theorem head
  {\newline}%                    space after theorem head (\newline forces line break)
  {}%                            theorem head spec

% Theorem-like environments
\theoremstyle{definition}
\newtheorem{definition}{Definition}[section]
\newtheorem{syntax}[definition]{Syntax}

% Use custom style for examples
\theoremstyle{examplestyle}
\newtheorem{example}[definition]{Example}

% Add extra space around examples
\BeforeBeginEnvironment{example}{\vspace{\baselineskip}}
\AfterEndEnvironment{example}{\vspace{0.5\baselineskip}}

% Fix verbatim indentation in examples
\BeforeBeginEnvironment{verbatim}{\par\vspace{0.5em}}
\AfterEndEnvironment{verbatim}{\vspace{0.5em}}

\theoremstyle{plain}
\newtheorem{theorem}[definition]{Theorem}
\newtheorem{proposition}[definition]{Proposition}
\newtheorem{lemma}[definition]{Lemma}
\newtheorem{corollary}[definition]{Corollary}

\theoremstyle{remark}
\newtheorem{remark}[definition]{Remark}
\newtheorem{note}[definition]{Note}
\newtheorem{compatibility}[definition]{R7RS Compatibility}

% Index entries
\newcommand{\indexentry}[1]{#1\index{#1}}
\newcommand{\indexmain}[1]{#1\index{#1|textbf}}

% Cross-references
\newcommand{\secref}[1]{Section~\ref{#1}}
\newcommand{\chapref}[1]{Chapter~\ref{#1}}
\newcommand{\figref}[1]{Figure~\ref{#1}}
\newcommand{\tabref}[1]{Table~\ref{#1}}
\newcommand{\appref}[1]{Appendix~\ref{#1}}

% Library references
\newcommand{\library}[1]{\texttt{(#1)}}
\newcommand{\proc}[1]{\texttt{#1}}
\newcommand{\var}[1]{\textit{#1}}

% Version information
\newcommand{\version}{0.2.0}
\newcommand{\baseversion}{R7RS-small}

% Title formatting - handled by algol60.cls

% Page layout adjustments
\setlength{\parskip}{0.5em}
\setlength{\parindent}{0pt}

% Table formatting
\newcolumntype{L}[1]{>{\raggedright\arraybackslash}p{#1}}
\newcolumntype{C}[1]{>{\centering\arraybackslash}p{#1}}
\newcolumntype{R}[1]{>{\raggedleft\arraybackslash}p{#1}}

% Float placement
\floatplacement{figure}{H}
\floatplacement{table}{H}

% Custom list environments
\newenvironment{grammarlist}{
    \begin{itemize}
    \setlength{\itemsep}{0pt}
    \setlength{\parskip}{0pt}
}{
    \end{itemize}
}

% Procedure specifications
\newenvironment{procspec}[1]{
    \vspace{1em}
    \noindent\textbf{Procedure:} \texttt{#1}
    \vspace{0.5em}
    \begin{quote}
}{
    \end{quote}
    \vspace{1em}
}

% Type specifications
\newenvironment{typespec}[1]{
    \vspace{1em}
    \noindent\textbf{Type:} \textsf{#1}
    \vspace{0.5em}
    \begin{quote}
}{
    \end{quote}
    \vspace{1em}
}

% Syntax specifications  
\newenvironment{syntaxspec}[1]{
    \vspace{1em}
    \noindent\textbf{Syntax:} \texttt{#1}
    \vspace{0.5em}
    \begin{quote}
}{
    \end{quote}
    \vspace{1em}
}

% Unicode-compatible code environment using listings package
% This supports Unicode characters and monospace fonts
\lstdefinestyle{lambdust-code-style}{
    language=Lisp,
    basicstyle=\ttfamily\footnotesize,
    commentstyle=\color{lambdust-comment}\itshape,
    keywordstyle=\color{lambdust-keyword}\bfseries,
    stringstyle=\color{lambdust-string},
    numberstyle=\color{lambdust-number},
    showstringspaces=false,
    breaklines=true,
    breakatwhitespace=true,
    tabsize=2,
    xleftmargin=0pt,
    xrightmargin=0pt,
    resetmargins=false,
    aboveskip=0.5em,
    belowskip=0.5em,
    captionpos=b,
    % Lambdust-specific keywords (same as before)
    morekeywords={define,lambda,let,let*,letrec,if,cond,case,and,or,not,
                  begin,do,when,unless,quote,quasiquote,unquote,syntax,
                  syntax-case,syntax-rules,with-syntax,datum,
                  Pi,Sigma,Refinement,the,perform,handle,lift,
                  spawn,send,receive,future,await,
                  foreign-call,foreign-data,capability,
                  prove,extract,qed,
                  Dynamic,Number,String,Boolean,Symbol,List,Vector,
                  Effect,Actor,Future,Capability},
    % Extended Unicode support
    extendedchars=true,
    inputencoding=utf8,
    literate={λ}{{\ensuremath{\lambda}}}1
             {→}{{\ensuremath{\rightarrow}}}1
             {∀}{{\ensuremath{\forall}}}1
             {∃}{{\ensuremath{\exists}}}1
             {Π}{{\ensuremath{\Pi}}}1
             {Σ}{{\ensuremath{\Sigma}}}1
             {⊢}{{\ensuremath{\vdash}}}1
             {≤}{{\ensuremath{\leq}}}1
             {≥}{{\ensuremath{\geq}}}1
             {≠}{{\ensuremath{\neq}}}1
             {∈}{{\ensuremath{\in}}}1
             {∉}{{\ensuremath{\notin}}}1
             {∅}{{\ensuremath{\emptyset}}}1
             {∪}{{\ensuremath{\cup}}}1
             {∩}{{\ensuremath{\cap}}}1
             {⊆}{{\ensuremath{\subseteq}}}1
             {⊇}{{\ensuremath{\supseteq}}}1
             {⟨}{{\ensuremath{\langle}}}1
             {⟩}{{\ensuremath{\rangle}}}1
}

% Define the lambdustcode environment
\lstnewenvironment{lambdustcode}{%
    \lstset{style=lambdust-code-style}%
}{}

% Add proper spacing hooks for lambdustcode environment
\BeforeBeginEnvironment{lambdustcode}{\par\vspace{0.5em}}
\AfterEndEnvironment{lambdustcode}{\vspace{0.5em}}

% Alternative inline code command 
\newcommand{\lambdustinline}[1]{\lstinline[style=lambdust-code-style]!#1!}

% Enhanced code command that supports Unicode
\renewcommand{\code}[1]{{\ttfamily\footnotesize #1}}